\section{Discussion: a decision table instead of a ranking}
\label{sec:discussion}

Advice in this area is usually organised by technology --- ``use Numba for this, use C for
that''. The measurements put the workload first and the technology second, for three reasons
that recur across all eight questions. The index that works is what the workload does, with the
technology and its price in the cells.

\paragraph{1. On compute-bound work the technologies are interchangeable; where allocation
dominates they are not.} On compute-bound scalar loops, Cython, Numba, C behind two different FFIs, and Rust
agree to within $\meas{pct:b1_compute-314:arraysum/c_ctypes|b1_compute-314:arraysum/numba}{0.2}\%$ on the reduction, and where they do not agree they differ by $\meas{pct:b1_compute-314:mandelbrot/rust_ctypes|b1_compute-314:mandelbrot/cython}{18}\%$ on
one kernel and $\meas{ratio:b1_compute-314:matmul/c_ctypes|b1_compute-314:matmul/numba}{1.8}\times$ on another, in both directions with no consistent winner --- against
the $32$--$\meas{ratio:b1_compute-314:matmul/cpython_loop|b1_compute-314:matmul/numba}{128}\times$ that compiling at all is worth (\S\ref{sec:rq1}). The measurements
therefore do not decide among them; what would --- build complexity, debuggability, code
ownership --- lies outside what this study measures. What \emph{does} distinguish outcomes is which class
the workload belongs to: the same four technologies span $52$--$61\times$ on the compute-bound
kernels and $7.3$--$13.6\times$ on the branchy and allocating ones (\S\ref{sec:rq2}) --- and
inside that second class they stop being interchangeable too, separating by $3.7$--$\meas{ratio:b2_branchy-314:binarytrees/numba|b2_branchy-314:binarytrees/rust_ctypes}{5.3}\times$
exactly where small objects are allocated.

\paragraph{2. The decisions that dominate are about what the program may do, not about which
tool does it.} Choosing
to write the hot loop in a language that can be compiled at all is worth more than any choice
made within that language --- except where the object model is the bound, and there the two are
the same size: compiling the binary trees at all buys Cython $\meas{ratio:b2_branchy-314:binarytrees/cpython_loop|b2_branchy-314:binarytrees/cython}{4.9}\times$, and choosing an
implementation that owns its own memory buys $3.5$--$\meas{ratio:b2_branchy-314:binarytrees/numba|b2_branchy-314:binarytrees/rust_ctypes}{5.3}\times$ on top of it
(\S\ref{sec:rq2}) --- and flattening
the data structure without compiling the loop made our graph traversal $\meas{pct:b2_branchy-314:bfs/cpython_flat|b2_branchy-314:bfs/cpython_loop}{12}\%$ \emph{slower}.
Deferring module execution is worth four orders of magnitude of import
latency (\S\ref{sec:rq4}) and changes when import errors surface. Allowing floating-point
reassociation is worth $\meas{ratio:b1_flags-314:arraysum/O3native|b1_flags-314:arraysum/O3native_ffast}{2.42}\times$ on a reduction, and it changes results --- and it needs no
compiler flag at all, since regrouping the accumulation in the source is worth $\meas{ratio:b1_flags-314:arraysum/accum_acc1|b1_flags-314:arraysum/accum_acc4}{2.74}\times$ at
unchanged flags (\S\ref{sec:rq1b}). Each is a decision about
what the program is allowed to do, and each dominates the tool choice made under it. The flags
below those decisions are not a substitute for them: moving from \code{-O2} to \code{-O3} was
worth nothing measurable on two of our three kernels and $\meas{pct:b1_flags-314:matmul/O2|b1_flags-314:matmul/O3}{0.1}\%$ on the third, while
\code{-march=native}, the one flag that did change the emitted code everywhere, made one kernel
$\meas{ratio:b1_flags-314:mandelbrot/O2|b1_flags-314:mandelbrot/O3native}{1.25}\times$ faster and another $\meas{ratio:b1_flags-314:matmul/O3native|b1_flags-314:matmul/O2}{1.80}\times$ \emph{slower} (\S\ref{sec:rq1b}).

\paragraph{3. Every gain has a cost, and the cost is the part that is usually not reported.}
JIT compilation costs $\meas{fact:b1_breakeven-314:facts.numba_compile_s:ms}{215}$\,ms of latency in the one case we measured, which is 7 calls or 7335
calls of work depending on input size (\S\ref{sec:rq1c}). Free-threading costs $1.05$--$\meas{ratio:b4_threads-ft:single_thread_arith/ft|b4_threads-gil:single_thread_arith/gil}{1.14}\times$ of
single-thread throughput (\S\ref{sec:rq5}). Static Python costs a rewrite in primitive types
and is $\meas{ratio:b6_static-fork_static:matmul96_int/fork_static/static_interp|b6_static-fork_static:matmul96_int/fork_static/boxed_python}{10.2}\times$ \emph{slower} if the JIT is not also enabled (\S\ref{sec:rq4}). An FFI costs
$59$--$\meas{med:b1_compute-314:call_overhead/c_ctypes:ns}{166}$\,ns per crossing, against $\meas{med:b1_compute-314:call_overhead/python_function:ns}{31}$\,ns for a pure-Python call: free for a $10\,\mu$s
kernel, a real cost per record (\S\ref{sec:rq2}).

\medskip
Table~\ref{tab:decision} is the result: the workload classes of \S\ref{sec:intro} with a
measured effect and a price in every row. It is the form we would put in front of an engineer.

\begin{table}[t]
  \centering
  \caption{Decision table from the measurements. Every factor is from the single measurement
    host of Table~\ref{tab:platform}; ``cost'' lists what the change also brings.}
  \label{tab:decision}
  \small
  \begin{tabular}{p{0.22\linewidth}p{0.25\linewidth}p{0.17\linewidth}p{0.26\linewidth}}
    \toprule
    workload & action that paid off & measured & cost / caveat \\
    \midrule
    numeric loop over an array
      & any compiled path (Cython / Numba / C / Rust)
      & $32\times$
      & a build step; the four routes are within $\meas{pct:b1_compute-314:arraysum/c_ctypes|b1_compute-314:arraysum/numba}{0.2}\%$ of each other here \\
    \dots{} and reassociation is acceptable
      & \code{fastmath} / \code{-ffast-math}
      & $\meas{ratio:b1_flags-314:arraysum/O3native|b1_flags-314:arraysum/O3native_ffast}{2.42}\times$ on top
      & results change ($4.6\cdot10^{-5}$ here) \\
    array operation a library already has
      & NumPy / vendor BLAS
      & up to $\meas{ratio:b1_compute-314:matmul/cpython_loop|b1_compute-314:matmul/numpy}{1495.9}\times$
      & only where the operation exists; $\meas{ratio:b1_compute-314:mandelbrot/cpython_loop|b1_compute-314:mandelbrot/numpy}{4.7}\times$ when it does not fit \\
    branch-heavy scan over flat data
      & Cython or Numba
      & $12\times$
      & C and Rust are only $8$--$\meas{pct:b2_branchy-314:tokenize/numba|b2_branchy-314:tokenize/c_ctypes}{17}\%$ further ahead \\
    allocation-heavy object graph
      & an allocator that is not CPython's: C or Rust behind an FFI, or Codon
      & $17.2$--$\meas{ratio:b2_branchy-314:binarytrees/cpython_loop|b2_branchy-314:binarytrees/rust_ctypes}{20.8}\times$ vs $3.9$--$\meas{ratio:b2_branchy-314:binarytrees/cpython_loop|b2_branchy-314:binarytrees/cython}{4.9}\times$
      & memory-ownership rewrite, or a recompile if Codon takes the source \\
    irregular graph traversal
      & compile it; a flat layout alone is not the win
      & $9.9\times$ (Cython and C within half a per cent)
      & flattening it in pure Python was worth $\meas{ratio:b2_branchy-314:bfs/cpython_loop|b2_branchy-314:bfs/cpython_flat}{0.89}\times$ \\
    object/attribute-heavy Python
      & upgrade to 3.11+
      & $2.0$--$\meas{ratio:b3_runtime-v310:attr_set/v310|b3_runtime-v314:attr_set/v314}{3.0}\times$ on attribute and method operations
      & 3.11 delivered most of it; 3.12 is slower than 3.11 on 14 of 16 operations; a dict
        literal on 3.14 runs at $\meas{ratio:b3_runtime-v310:create_dict/v310|b3_runtime-v314:create_dict/v314}{0.74}\times$ its 3.10 rate \\
    thread-parallel CPU-bound Python
      & free-threaded build
      & $5.6$--$\meas{ratio:b4_threads-ft:py_branchy/ft/T1_control|b4_threads-ft:py_branchy/ft/T8}{6.1}\times$ on 8 threads
      & $1.05$--$\meas{ratio:b4_threads-ft:single_thread_arith/ft|b4_threads-gil:single_thread_arith/gil}{1.14}\times$ single-thread tax, paid by the build not the lock \\
    thread-parallel work already in C
      & threads on any build
      & $5.8$--$\meas{ratio:b4_threads-gil:c_ffi_uniform/gil/T1|b4_threads-gil:c_ffi_uniform/gil/T8}{7.2}\times$ on 8 threads, balanced kernel
      & the GIL is irrelevant here; the imbalanced kernel reaches $1.8$--$\meas{ratio:b4_threads-gil:c_ffi_mandelbrot/gil/T1_control|b4_threads-gil:c_ffi_mandelbrot/gil/T8}{3.1}\times$ \\
    interpreter build tuning
      & PGO with full LTO, not PGO alone
      & $\meas{geomean:b5_buildcfg-plain|b5_buildcfg-pgo_ltofull}{1.10}\times$ ($\meas{geomean:b5_buildcfg-plain|b5_buildcfg-pgo_ltofull_native}{1.12}\times$ adding \code{-march=native})
      & PGO alone is $\meas{geomean:b5_buildcfg-plain|b5_buildcfg-pgo}{1.02}\times$; attribute writes lose in all six, $0.88$--$\meas{ratio:b5_buildcfg-plain:attr_set/plain|b5_buildcfg-pgo:attr_set/pgo}{0.98}\times$ \\
    CPython's own tier-2 JIT
      & measure it: it helps more operations than it hurts
      & $\meas{geomean:b5_buildcfg-plain|b5_buildcfg-jit}{1.04}\times$ overall
      & 6 of 16 operations slower (integer arithmetic and attribute writes $0.88$--$\meas{ratio:b5_buildcfg-plain:int_arith/plain|b5_buildcfg-jit:int_arith/jit}{0.89}\times$);
        still labelled experimental \\
    pure Python that must stay pure Python
      & PyPy
      & $3.2$--$\meas{ratio:b1_compute-314:matmul/cpython_loop|b1_compute-pypy311:matmul/cpython_loop}{42.8}\times$
      & a \code{ctypes} call passing a $2$\,MB buffer costs $\meas{ratio:b2_branchy-pypy311:tokenize/c_ctypes|b2_branchy-314:tokenize/c_ctypes}{5.3}\times$ more than CPython's;
        no Numba; ecosystem support is thinning \\
    a reduction bound by its dependency chain
      & regroup the accumulation in the source
      & $\meas{ratio:b1_flags-314:arraysum/accum_acc1|b1_flags-314:arraysum/accum_acc4}{2.74}\times$ at unchanged flags, above \code{-ffast-math}'s $\meas{ratio:b1_flags-314:arraysum/O3native|b1_flags-314:arraysum/O3native_ffast}{2.42}\times$
      & a different summation order, whether a flag or the source grants it
        (\S\ref{sec:rq1b}) \\
    typed hot code, willing to rewrite
      & CinderX Static Python + JIT
      & $\meas{ratio:b6_static-fork_static_adaptive:matmul96_int/fork_static_adaptive/boxed_jit|b6_static-fork_static_adaptive:matmul96_int/fork_static_adaptive/static_jit}{11.4}\times$ on top of the JIT
      & boot sequence documented only in a tutorial; $\meas{ratio:b6_static-fork_static:matmul96_int/fork_static/static_interp|b6_static-fork_static:matmul96_int/fork_static/boxed_python}{10.2}\times$ slower without the JIT, $\meas{ratio:b6_static-fork_static_adaptive:matmul96_int/fork_static_adaptive/static_interp|b6_static-fork_static_adaptive:matmul96_int/fork_static_adaptive/boxed_python}{5.8}\times$ slower
        even with the specialising interpreter path \\
    startup dominated by imports
      & Cinder lazy imports
      & $\meas{med:b6_cinderx-fork:import_200_modules_eager/fork:ms}{19.7}$\,ms $\to$ under $1\,\mu$s
      & fork-only; defers import errors to first use; first touch costs $\meas{med:b6_cinderx-fork:import_200_modules_first_touch/fork:ms}{0.35}$\,ms \\
    mixed pipeline
      & attack the largest stage, and stop there
      & $\meas{ratio:b7_pipeline-gil:end_to_end/gil/v0_pure|b7_pipeline-gil:end_to_end/gil/v2_numba_native}{1.35}\times$ from one stage
      & Numba on a \meas{share:b7_pipeline-gil:stage_aggregate/gil/v0_pure|b7_pipeline-gil:end_to_end/gil/v0_pure}{3.6}\% stage gave $\meas{ratio:b7_pipeline-gil:end_to_end/gil/v0_pure|b7_pipeline-gil:end_to_end/gil/v1_numba}{1.01}\times$; adding threads gave $\meas{ratio:b7_pipeline-gil:end_to_end/gil/v0_pure|b7_pipeline-gil:end_to_end/gil/v3_native_threads}{0.72}\times$ (GIL) and
        $\meas{ratio:b7_pipeline-ft:end_to_end/ft/v0_pure|b7_pipeline-ft:end_to_end/ft/v3_native_threads}{1.28}\times$ (free-threaded) \\
    \bottomrule
  \end{tabular}
\end{table}

\paragraph{Four questions to ask of a performance figure.} Working through the eight research
questions of this paper produced a short checklist. What was the input size --- a break-even point moves by three
orders of magnitude across ours (\S\ref{sec:rq1c})? Was the baseline the same program --- a
vendor BLAS against a naive triple loop is not a language comparison (\S\ref{sec:rq1})? Was one
variable changed, or two --- an interpreter's build options are worth several releases of
interpreter work, so a version sweep on differently-built binaries measures neither
(\S\ref{sec:rq3}, \S\ref{sec:rq6})? And what did it cost --- on the code it does not
accelerate, at the boundary it introduces, in the semantics it relaxes (\S\ref{sec:rq5},
\S\ref{sec:rq2}, \S\ref{sec:rq1b})? A figure that survives those four is worth acting on.
